\documentclass[11pt]{article}
\usepackage{amsmath}
\usepackage{geometry}
\geometry{a4paper, margin=1in}

\begin{document}

\section*{Shape Equations for Planar Membrane Wrap}

Based on Agudo-Canalejo (2021) "Particle engulfment by strongly asymmetric membranes with area reservoirs", the explicit axisymmetric shape equations for the unbound membrane segment are given below. 

To facilitate numerical integration over an unknown physical arc length, we introduce a scaling parameter $\Omega$ representing the total arc length of the unbound membrane. The integration variable $s$ is rescaled to a fixed dimensionless domain $s \in [0, 1]$. The physical derivatives are scaled by $\Omega$ via the chain rule. We also augment the system with the height coordinate $z$.

The derivatives with respect to the rescaled arc length $s$ are:

\begin{align}
\dot{\psi} &= \Omega u \\
\dot{u} &= \Omega \left( -\frac{\cos \psi}{x} u + \frac{\sin \psi \cos \psi}{x^2} + \frac{\sin \psi}{2 \cos \psi } \left[ \frac{\sigma}{\kappa} - u^2 + \frac{\sin^2 \psi}{x^2} + \frac{4 m \sin \psi}{x} \right] \right) \\
\dot{x} &= \Omega \cos \psi \\
\dot{z} &= -\Omega \sin \psi
\end{align}

where $\sigma$ is the effective tension combining the mechanical tension $\Sigma$ and the spontaneous tension $2 \kappa m^2$:

\begin{equation}
\sigma = \Sigma + 2 \kappa m^2
\end{equation}

\paragraph{Boundary Conditions and Shooting Method}
The initial conditions at $s=0$ correspond to a smooth matching with the boundary of the particle. The wrapping angle is $\phi$.
\begin{align*}
\psi(0) &= \phi \\
u(0) &= u_0 \quad (\text{guessed parameter}) \\
x(0) &= R_{pa} \sin \phi \\
z(0) &= R_{pa} (1 - \cos \phi) \quad (\text{or } 0 \text{ depending on reference})
\end{align*}

We integrate from $s=0$ to $s=1$. The target boundary conditions for a flat far-field membrane are:
\begin{align*}
\psi(1) &= 0 \\
u(1) &= 0
\end{align*}

Because we have two target conditions at the right boundary, we require two shooting parameters: the initial curvature $u_0$ and the total arc length $\Omega$. 

\paragraph{Numerical Tips \& Convergence Fallback}
\begin{itemize}
    \item \textbf{Small $x$ Instability:} Check if $x$ is too small (e.g., $< 0.1$), which could lead to numerical instabilities. If $R_{pa}$ is large enough, it should not be an issue. Alternatively, expand the shape equation similarly to the poles in the spherical vesicle case.
    \item \textbf{Convergence Fallback:} The membrane asymptotically approaches the flat state over a characteristic length scale $\lambda = \sqrt{\kappa/\sigma}$. Forcing exactly $\psi=0$ and $u=0$ at a finite $\Omega$ can sometimes cause the root-finding algorithm to fail. \textit{If convergence is not obtained}, drop $\Omega$ as a shooting parameter. Instead, fix $\Omega$ to a large constant value (e.g., $\Omega \approx 10\lambda$) and perform a 1D shooting method optimizing only $u_0$ to minimize $\psi(1)^2 + u(1)^2$.
\end{itemize}

\paragraph{Energy of Unbound Membrane}
The shape equations (1)-(4) are augmented with an equation for the membrane free energy (calculated with respect to the energy of a flat membrane reference state):

\begin{equation}
\dot{F}_{\text{me}} = \Omega \pi \kappa x \left( \left( \frac{\dot{\psi}}{\Omega} + \frac{\sin \psi}{x} \right)  \left( \frac{\dot{\psi}}{\Omega} +  \frac{\sin \psi}{x} + 4m \right) + \frac{2 \sigma}{\kappa} (1 - \cos \psi) \right)
\end{equation}

Integrating (6) from $s = 0$ to $s = 1$ yields the free energy of the unbound membrane segment $F_{me}^{un}$ for any given wrapping angle $\phi$.

\paragraph{Total Energy of the Process}
The bound segment of the membrane follows the shape of the particle (a spherical cap). Its free energy, relative to a flat state, is:
\begin{equation}
    F_{me}^{bo} = 4 \pi \kappa (1 + 2mR_{pa})(1 - \cos \phi)  + \sigma \pi R_{pa}^2(1-\cos \phi)^2
\end{equation}

The adhesion free energy is:
\begin{equation}
    F_{ad}(\phi) = - |W| 2 \pi R_{pa}^2 (1 - \cos \phi)
\end{equation}

The total free energy of the system as a function of the wrapping angle $\phi$, defining the energy landscape for engulfment, is given by:
\begin{equation}
    F_{tot}(\phi) = F_{me}^{un} + F_{ad}(\phi) + F_{me}^{bo}
\end{equation}

\end{document}